%% LyX 2.4.4 created this file.  For more info, see https://www.lyx.org/.
%% Do not edit unless you really know what you are doing.
\documentclass[english]{article}
\usepackage[T1]{fontenc}
\usepackage[latin9]{inputenc}
\usepackage{enumitem}
\usepackage{amsmath}
\usepackage{amsthm}
\usepackage{cancel}
\usepackage{geometry}
\geometry{verbose,tmargin=1in,bmargin=1in,lmargin=1in,rmargin=1in}

\makeatletter
%%%%%%%%%%%%%%%%%%%%%%%%%%%%%% Textclass specific LaTeX commands.
\numberwithin{equation}{section}
\numberwithin{figure}{section}
\newlength{\lyxlabelwidth}      % auxiliary length 

%%%%%%%%%%%%%%%%%%%%%%%%%%%%%% User specified LaTeX commands.
\usepackage{fancyhdr}
\usepackage{lastpage}
\pagestyle{fancy}
\usepackage{enumitem}
\usepackage{ifthen}
%sets math fonts to be more readable
\everymath=\expandafter{\the\everymath\displaystyle}
%\DeclareMathSizes{10}{10}{10}{10}
\usepackage{multicol}

\usepackage{tikz}
\usepackage{tikz-3dplot}
\usetikzlibrary{calc,arrows}

\usetikzlibrary{plotmarks}
\usepackage{pgfplots}
\pgfplotsset{compat=newest}

\usepackage{calc}
\usepackage[nomessages]{fp}

\usepackage{advdate}    % Advancing/saving dates
\usepackage{datetime}   % Dates formatting
\usepackage{datenumber} % Counters for dates
% Added by lyx2lyx
\setlength{\parskip}{\smallskipamount}
\setlength{\parindent}{0pt}

\makeatother

\usepackage{babel}
\begin{document}
\renewcommand{\author}{Dr.\,James Ashe}

\newcommand{\classNum}{331}

\newcommand{\classSec}{A}

\newcommand{\nameType}{Name:}

\newcommand{\examType}{Worksheet 12.3}

\renewcommand{\date}{\formatdate{16}{10}{2013}}

%%First page's header%%%%%%%%%%%%%%%%%%%%%%
\fancypagestyle{firststyle} %%header settings for first pagr
{
	\fancyhead[L]{MTH \classNum{}(\classSec{}) \author{}\\
	\textbf{\examType{}} \date{}}
	\fancyhead[C]{\textbf{\nameType}}
	\setlength{\headsep}{14pt}
}
%%%%%%%%%%%%%%%%%%%%%%%%%%%%%%%%%%%%%%%%%%

%%Next page's header%%%%%%%%%%%%%%%%%%%%%%%
\setlist{nolistsep}
\lhead{\classNum{}(\classSec{})} 
\chead{\textbf{\nameType}} 
\cfoot{\thepage{}~of~\pageref{LastPage}} 
\setlength{\headsep}{5pt} 
%%%%%%%%%%%%%%%%%%%%%%%%%%%%%%%%%%%%%%%%%%%

%%Various settings; see the preamble for other settings%%%%%
%%\setenumerate[0]{leftmargin=12pt,itemindent=0pt,labelwidth=0pt}
\renewcommand{\labelenumi}{\textbf{\arabic{enumi}.}}
\renewcommand{\labelenumii}{\textbf{\alph{enumii}.}}
\thispagestyle{firststyle}
%%%%%%%%%%%%%%%%%%%%%%%%%%%%%%%%%%%%%%%%%%
\newcommand{\xmax}{10}
\newcommand{\xmin}{-10}
\newcommand{\xstep}{0} %must be xmin+n
\newcommand{\ymax}{10}
\newcommand{\ymin}{-10}
\newcommand{\ystep}{0} %must be ymin+n

\newcommand{\xspace}{\vspace{1.2cm}}
\begin{enumerate}[resume]
\item Let $\mathbf{v=}\left\langle a,b,c\right\rangle $ and let $\alpha$,
$\beta$, and $\gamma$ be the angles between $\mathbf{v}$ and the
positive $x$-axis, the positive $y$-axis, and the positive $z$-axis,
respectively (see the figure below).

\begin{enumerate}[resume]
\item Prove that $\cos^{2}\alpha+\cos^{2}\beta+\cos^{2}\gamma=1$.
\item Find a vector that makes a $45^{\circ}$ angle with $\mathbf{i}$
and $\mathbf{j}$. What angle does it make with $\mathbf{k}$?
\item Is there a vector that makes a $30^{\circ}$ angle with $\mathbf{i}$
and $\mathbf{j}$? Explain.
\item Find a vector $\mathbf{v}$ such that $\alpha=\beta=\gamma$. What
is the angle?
\end{enumerate}
\tdplotsetmaincoords{70}{130}% %%specifes orientation
%%\tdplotsetrotatedcoords{0}{20}{0}% %%specifes angle of rotation
\begin{tikzpicture} 		
	[tdplot_main_coords,
		scale=1,
		cube/.style={purple}, 
		grid/.style={very thin,gray}, 
		axis/.style={->,black,thick},
		vector/.style={very thick,red,-stealth},
		vector guide/.style={dashed,red,thick},
		angle/.style={black,thick}]
	
	%draw a grid in the x-y plane 	
	\foreach \x in {-0.25,0,...,1.25} 		
	\foreach \y in {-0.25,0,...,1.25} 		
	{ 			
		%\draw[grid] (\x,-0.25) -- (\x,1.25); 			
		%\draw[grid] (-0.25,\y) -- (1.25,\y); 		
	} 		
	
	
	
	%draw the axes 	
	\draw[axis] (-1,0,0) -- (3.5,0,0) node[anchor=south]{$x$}; 	
	\draw[axis] (0,-1,0) -- (0,3.5,0) node[anchor=west]{$y$}; 	
	\draw[axis] (0,0,-1) -- (0,0,3.5) node[anchor=west]{$z$};
	
	%draw the top and bottom of the cube 	
	%\draw[cube] (0,0,0) -- (0,1,0) -- (1,1,0) -- (1,0,0) -- cycle; 	
	%\draw[cube] (0,0,1) -- (0,1,1) -- (1,1,1) -- (1,0,1) -- cycle; 
	
	%draw the edges of the cube 	
	%\draw[cube] (0,0,0) -- (0,0,1); 	
	%\draw[cube] (0,1,0) -- (0,1,1); 	
	%\draw[cube] (1,0,0) -- (1,0,1); 	
	%\draw[cube] (1,1,0) -- (1,1,1);

	%standard tikz coordinate definition using x, y, z coords 	
	\coordinate (O) at (0,0,0);
	%tikz-3dplot coordinate definition using r, theta, phi coords 	
	\tdplotsetcoord{P}{.8}{55}{60}

	%draw vector
	%\draw[vector] (O) -- (P);
	%\draw[vector guide] (O) -- (Pxy); \draw[vector guide] (Pxy) -- (P);
	\tdplotsetthetaplanecoords{90} %rotates around z-axis
	\tdplotdrawarc[tdplot_rotated_coords,angle,red]{(O)}{2.5}{0}{33.6}{anchor=west}{$\gamma$}; %{center}{radius}{start}{end angle}

	\tdplotsetthetaplanecoords{90}	
	\tdplotdrawarc[tdplot_rotated_coords,angle,purple]{(O)}{2}{33.69}{90}{anchor=west}{$\beta$}; %{center}{radius}{start}{end angle}

	%\tdplotsetcoord{P}{0}{45}{45};	
	\tdplotsetrotatedthetaplanecoords{33.6}
	\tdplotdrawarc[tdplot_rotated_coords,angle,blue]{(O)}{1.5}{180}{90}{anchor=south}{$\alpha$}; %{center}{radius}{start}{end angle}
	
	\draw[->,very thick,black,-stealth] (0,0,0) -- (0,2,3);
	%draw an arc illustrating the angle defining the orientation 

\end{tikzpicture}
\end{enumerate}

\end{document}
